\chapter{概率与统计}
\setcounter{choicecounter}{0}

%% ---- 知识框架 ---- %%
\begin{knowledgeframe}
\begin{itemize}
	\item \textbf{相关系数与决定系数}：$|r|$ 越大线性相关性越强；$R^2$ 越大拟合效果越好。
	\item \textbf{$\chi^2$ 独立性检验}：适用于抽样数据，不适用于普查数据。
	\item \textbf{正态分布}：$\xi \sim N(\mu, \sigma^2)$ 的对称性——$P(\xi \ge \mu+a)=P(\xi \le \mu-a)$。
	\item \textbf{独立事件的概率}：各项目独立时，总概率为各项目概率的乘积。
	\item \textbf{分布列与期望}：离散型随机变量 $X$ 的 $E(X)=\sum x_i p_i$。
\end{itemize}
\end{knowledgeframe}

%% ---- 第一节：统计基础 ---- %%
\section{统计基础}

\startexercise

\subsection*{A组}

\begin{choice}{2025雅礼月考八（多选）}{统计概念辨析}
	下列说法正确的是
	\begin{tasks}(2)
		\task 线性相关系数 $r$ 越小，两个变量的线性相关性越弱
		\task 在线性回归模型中，决定系数 $R^2$ 越大，表示残差平方和越小，即模型的拟合效果越好
		\task $\chi^2$ 独立性检验方法不适用于普查数据
		\task 已知随机变量 $\xi \sim N(0, \sigma^2)$，若 $P(\xi \ge -2) = 0.8$，则 $P(\xi > 2) = 0.2$
	\end{tasks}
\end{choice}

\stopexercise

%% ---- 第二节：概率与分布 ---- %%
\section{概率与分布}

\startexercise

\subsection*{B组}

\begin{exercise}[2022全国甲理19]
	甲、乙两个学校进行体育比赛，比赛共设三个项目，每个项目胜方得 10 分，负方得 0 分，没有平局。三个项目比赛结束后，总得分高的学校获得冠军。已知甲学校在三个项目中获胜的概率分别为 $0.5, 0.4, 0.8$。各项目的比赛结果相互独立。
	\begin{enumerate}
		\item[(1)] 求甲学校获得冠军的概率；
		\item[(2)] 用 $X$ 表示乙学校的总得分，求 $X$ 的分布列与期望。
	\end{enumerate}
\end{exercise}

\stopexercise
